% =========================================================================
% Section 2 — Related Work (Korean translation)
% File: sections/2_related_work.tex  (\input fragment; no \documentclass)
% Model name: DefectSegAM (exact spelling, everywhere)
% =========================================================================

\section{관련 연구}
\label{sec:related}

\subsection{레이어 단위 분말 베드 모니터링과 DSCNN--Peregrine 계열}
\label{sec:related-monitoring}

초기의 in-situ 모니터링은 결함 검출을 공정 수준의 신호 분류 문제로 다루었다.
Okaro 등은 49개 인장 바(tensile bar)에서 획득한 광다이오드 신호로
준지도(semi-supervised) 모델을 학습하여 약 77\%의 성공률로 불량 공정을
판별하였고~\cite{Okaro2019SemiSup}, Baumgartl 등은 H13 공구강의 off-axis
적외선 서모그램을 경량 합성곱 네트워크로 분류하여 평균 균형 정확도(mean
balanced accuracy) 96.8\%를 달성하였다~\cite{Baumgartl2020IRCNN}. 이러한
분류기는 적시에 경보를 제공하지만 결함을 국소화하지는 못한다. 공간 분해능은
레이어 단위 분말 베드 촬영 및 용융풀(melt-pool) 이미징과 함께 확보되기
시작하였다: Scime와 Beuth는 엔지니어링된 컴퓨터 비전 피처로 여섯 가지 분말
베드 이상 유형을 분류하였고~\cite{Scime2018Anomaly}, Inconel~718에 대해
용융풀 형태를 키홀 및 balling 영역과 연관지었으며~\cite{Scime2019MeltPool},
Gobert 등은 다중 모달 통합 후 약 85\%의 정확도로 서포트 벡터 머신을 이용해
고해상도 레이어 이미지의 불연속(discontinuity)을
검출하였다~\cite{Gobert2018InProcess}.

Scime 등의 DSCNN~\cite{Scime2020DSCNN}은 이 연구 계열을 여덟 개 클래스
온톨로지를 갖는 단일한 장비 불특정(machine-agnostic) 픽셀 단위 분말 베드
세그멘테이션 모델로 통합하였다 \TODO{투고 전에 원저 DSCNN의 학습 데이터
규모(픽셀, 레이어, 빌드 수)를 2020년 논문과 대조하여 검증할 것}.
Peregrine 모니터링 플랫폼에 배포되어 ORNL Co-Registered 데이터셋(Peregrine
v2023-11)~\cite{ORNL_CoRegistered_Dataset}과 함께 제공되는 양산 후속 모델은
온톨로지를 12개 클래스로 확장하고(recoater hopping, incomplete spreading,
misprint, under-melting 추가) 모든 레이어에 대해 $1842\times1842$ 크기의
클래스별 마스크를 생성한다. 이 온톨로지의 클래스는 관찰되는 이미지에
따라 개략적으로 갓 도포된 분말 베드의 이상(recoater hopping, recoater
streaking, incomplete spreading, debris), 용융된 표면의 이상(spatter,
misprint, 키홀 전조로서의 over-melting, lack-of-fusion 전조로서의
under-melting), 그리고 두 이미지 모두에서 관찰되는 이상(swelling,
super-elevation)으로 나눌 수 있다. 이 구분은 저자들의 정리이며, 외형
기준으로 다시 정리한 결과는 Table~\ref{tab:taxonomy}에 제시한다. 함께
공개된 Peregrine 데이터셋 릴리스는 전문가 어노테이션을 공개하였으나 그
양은 적다: v2021-03 릴리스에서는
기술(technology)당 20개의 어노테이션 레이어가~\cite{Peregrine_v2021_03},
v2022-10 릴리스에서는 여섯 개 장비--소재 조합 각각에 대해 5개에서 26개의
전문가 어노테이션 학습 레이어가 제공된다~\cite{Peregrine_v2022_10}. 이처럼
DSCNN--Peregrine 계열은 픽셀 단위 세그멘테이션을 플랫폼 규모 분말 베드
모니터링의 근간으로 확립하였으나, 이는 수십 개의 라벨링된 레이어에 대한
장비별 지도학습이라는 비용을 수반한다.

이 계열이 물려준 결함 어휘 자체도 검토가 필요하다. 예컨대 금속 PBF 결함
리뷰는 결함을 발생 원인과 함께 정의하는 것이 관례이며, Grasso와 Colosimo의
리뷰는 기공, balling, 기하 변형, 균열과 박리 등의 결함 유형을 각각의 발생
원인과 연결하여 정리한다~\cite{Grasso2017Review}. 반면 분말 베드 이미지의
이상 분류는 고해상도 레이어 이미징으로 관찰 가능한 공정 오류를 목록화한
초기 연구~\cite{zurJacobsmuhlen2013HighRes}와 Scime와
Beuth~\cite{Scime2018Anomaly} 이래 시각 형상을 기준으로 클래스를 정의해
왔으나, 그 명칭과 설명에는 여전히 원인이 섞여 있다. ASTM E3353-22 표준
지침~\cite{ASTME3353}도 공개된 4절(Significance and Use)에서 L-PBF의
광학·열 in-process 모니터링으로 검출할 분말층 이상의 원인을 리코팅
오류(skipping, scraping, insufficient powder delivery, part strike)와 파트
형성 오류(distortion, humping, balling, superelevation)로 범주화하고,
레이어 이미지에서의 상세 기술은 본 연구가 검토하지 못한 8절(Layer
Imaging)로 미룬다. 같은 4절은 과열(overheating)·overmelting 항목을 따로
두어 그 징후로 과열 영역의 변색(scorching)과 humping 등의 표면 변화를
들고, 과도한
spatter/ejecta를 그 하위 항목으로 기술한다. 이를 DSCNN--Peregrine 어휘와
대응시키면 skipping과 insufficient powder delivery는 모두 Incomplete
Spreading에, scraping은 Recoater Streaking에, part strike는 Recoater
Hopping에 대응한다. Distortion, humping, superelevation은 Swelling 또는
Super-Elevation에(humping의 대응은 잠정적이다; Table~\ref{tab:taxonomy}),
overmelting은 Over Melting에, spatter/ejecta는 Spatter에 대응하며,
balling에는 L-PBF 클래스 공간에 대응 클래스가 없다(Peregrine v2021-03의
전자빔 PBF 소스에만 Balling 클래스가 정의되어
있다~\cite{Peregrine_v2021_03}). 이 범주는 원인 기준이며, 본 연구가 검토한
공개부는 skipping과 insufficient powder delivery의 레이어 이미지 외형
차이를 제시하지 않는다. 또한 저자들이 검토한 ORNL
코퍼스와 사람 어노테이션 코퍼스도 둘을 구분하여 어노테이션하지
않는다(Section~\ref{sec:method-taxonomy}). 반대로 Incomplete Spreading의
정의는 분말 도포량 부족이라는 원인만 규정하므로, 얕은 미도포 자국과 분말이
무너져 내린 형상처럼 외형이 크게 다를 수 있는 이상이 한 용어에 포함될 수
있다. Peregrine 어휘 역시 외형 기준 클래스에 에너지 밀도라는 원인 기준의
Over Melting과 Under Melting, CAD 의도 기준의 Misprint를 섞어 쓴다. 또한
Peregrine v2021-03과 v2022-10 릴리스의 정의는 Swelling과 Super-Elevation을
모두 파트가 분말 베드 위로 솟은 외형(warps up, extends above the powder
bed)으로 기술하면서 국소 에너지 과잉과 분말 피복 부족·잔류 응력이라는
원인 절로만 구분하므로~\cite{Peregrine_v2021_03,Peregrine_v2022_10}, 픽셀
단위 어노테이션의 판정 기준이 클래스마다 다르다. ORNL Co-Registered
데이터셋(v2023-11)의 정의는 두 클래스를 이와 다르게 기술하므로 릴리스
사이에서도 판정 기준이 일치하지 않는다(Section~\ref{sec:method-taxonomy}).
이러한 원인 기준과 외형 기준의 혼재가 Section~\ref{sec:method-taxonomy}에서
외형 기준 결함 분류 체계를 제안하는 동기가 된다.

\subsection{Ex-situ ground truth로 지도되는 결함 검출}
\label{sec:related-exsitu}

이와 병행하는 연구 계열은 사람 어노테이션을 ex-situ 계측, 주로 XCT로
대체한다. Snow 등은 XCT에서 유도한 결함 라벨로 레이어 단위 영상에 대해
신경망을 학습하여, 학습 빌드 내에서는 200~$\mu$m보다 큰 결함에 대해 93.5\%의
정확도를 달성하였으나 미학습 빌드에서는 87.3\%에 그쳤으며, 이는 빌드 간
일반화 성능 저하를 조기에 정량화한 결과이다~\cite{Snow2021Layerwise}. Wong
등은 U-Net 계열 모델로 AM XCT 슬라이스의 기공과 균열을 분할하여 도메인 내
최고 mean intersection-over-union 0.993을 달성하였다~\cite{Wong2022UNet}.
영상 이외의 정보로는, Estalaki 등이 복셀 수준 열 특징으로부터 SS~304L의
micro-CT 라벨 기공률을 랜덤 포레스트로 예측하였고(테스트 F1 0.96
이상)~\cite{Estalaki2022Voxel}, Smoqi 등은 이중 파장 파이로미터로부터 물리
기반 용융풀 시그니처를 도출하여 기공의 유형과 심각도를 F1 0.95 이상으로
분류하였다~\cite{Smoqi2022Meltpool}. 최근 연구는 센싱 스택을 다변화하고
있다: Ero 등은 근적외선 광학 단층촬영(optical tomography)을
자기조직화지도(self-organizing map) 및 U-Net과 결합하여 EOS M290에서
lack-of-fusion 및 키홀 기공을 검출하였고~\cite{Ero2023OT}, Kumar 등은
열화상과 CCD 영상을 설명 가능한 네트워크로 융합하여 정확도를 95\%(열화상
단독)와 98.28\%(CCD 단독)에서 다중 모달 입력 시 99.72\%까지
향상시켰다~\cite{Kumar2024MultiSource}. 이들 연구는 도메인 내에서는 우수한
성능을 달성하지만, 라벨 생산 비용이 높고, 모델이 장비별·소재별로 학습되며,
Snow 등이 관찰한 바와 같이 표면에서 관측 가능한 단서가 없는 결함은 레이어
영상만으로는 여전히 검출하기 어렵다~\cite{Snow2021Layerwise}.

\subsection{AM 분야의 파운데이션 모델과 라벨 효율적 세그멘테이션}
\label{sec:related-foundation}

주류 컴퓨터 비전에서는 자기지도 파운데이션 모델(foundation model)이 태스크별
피처를 처음부터 학습하는 방식의 대안으로 자리 잡고 있다. DINOv2는 인코더의
미세조정 없이도 세그멘테이션으로 전이되는 범용 밀집 피처를
생성하고~\cite{Oquab2024DINOv2}, 레지스터 토큰은 비전 트랜스포머의 밀집 예측
성능을 저해하는 피처 맵 아티팩트를 제거하며~\cite{Darcet2024Registers},
dense prediction transformer (DPT)는 여러 트랜스포머 깊이에서 추출한 피처를
전체 해상도 예측으로 재조립하는 방법을 확립하였다~\cite{Ranftl2021DPT}. AM
분야에서 이 패러다임은 이제 막 나타나기 시작한 단계이다. Era 등은 프롬프트
기반(promptable) Segment Anything Model을 L-PBF CoCr XCT 스캔의 기공
세그멘테이션에 적용하여 태스크별 라벨 없이 Dice 유사도 계수 80\%를
상회하였고~\cite{Era2025SAM}, Luo 등은 광다이오드 유래 비트맵으로부터의 물성
예측에 ImageNet 사전학습 ResNet-101 전이학습을
사용하였다~\cite{Luo2025TransferTi}. 전이학습과 도메인 적응은 분포 시프트에
대한 고전적 대응이지만~\cite{Pan2010TransferLearning,Quinonero2009DatasetShift,Torralba2011DatasetBias},
저자들이 아는 한 고정된 자기지도 파운데이션 인코더를 다중 장비 분말 베드 결함
세그멘테이션---장비, 광학계, 소재 간의 시프트가 곧 장비별 지도학습
네트워크의 실패 모드가 되는 설정---에 적용한 선행 연구는 없다.

\subsection{결함 맵에서 기계적 물성으로}
\label{sec:related-property}

금속 AM에서 물성 예측은 오랫동안 명목 공정 파라미터로부터의 회귀를
의미하였으며, 이러한 접근은 동일한 설정 하에서 발생하는 부품 간·부품 내
변동을 설명하지
못한다~\cite{DebRoy2018PMS,DebRoy2021NatRevMater,Mahmoud2021LPBFreview}.
In-situ 데이터는 이 한계에 대응하지만 파괴 시험 라벨이 희소하므로, 대부분의
연구는 수십에서 수백 개 수준의 시편을 대상으로 수행된다. Xie 등은 적외선 열
이력으로부터 directed energy deposition으로 제작한 IN718의 인장강도를
ResNet-18로 예측하였고(135개 쿠폰에서 $R^2 \approx
0.7$)~\cite{Xie2021UTSResNet}; Fang 등은 시뮬레이션된 열 이력에 대한 1차원
합성곱 네트워크로 학습 데이터에서 $R^2 = 0.96$, 테스트 데이터에서 0.67에
도달하였으며 이는 단일 냉각 속도 지표의 0.32와
대비된다~\cite{Fang2022Inconel1DCNN}; Moon 등은 CT 기공 통계량과 표면
조도로부터 L-PBF Ti-6Al-4V의 피로 수명을 예측하였고(Pearson 상관계수
$\approx 0.94$)~\cite{Moon2021Fatigue}; Luo 등은 전이학습으로 UTS $R^2 =
0.89$와 연신율 $R^2 = 0.96$을 보고하였다~\cite{Luo2025TransferTi}. Scime
등의 voxel-wise property prediction model (VPPM)
프레임워크~\cite{VPPM_Baseline}는 규모 면에서 예외적이다: 약 230~GB
(${\sim}230$~GB) 규모 공개 데이터셋~\cite{ORNL_CoRegistered_Dataset}의 다섯
개 빌드에 걸친 6{,}299개의 SS~316L 인장 시편을 대상으로, DSCNN 결함
마스크에서 유도된 채널을 포함한 21개의 핸드크래프트 피처를 $1 \times 1
\times 3.5$~mm 슈퍼복셀 단위로 집계하여 물성별 다층 퍼셉트론에 입력하고,
시편 단위 5-fold 교차검증 하에서 in-situ 정보가 없는 베이스라인 대비 UTS
예측 오차의 61\% 감소를 보고하였다. 이 프레임워크는 결함 세그멘테이션을
물성 예측의 입력으로 전환함으로써 세그멘테이션 품질을 하류 유용성으로 측정
가능하게 만들며, 본 연구는 Section~\ref{sec:results}에서 이 지렛대를
활용한다.

\paragraph{연구 공백과 본 연구의 위치}
이러한 연구 계열 전반에서 세 가지 결핍이 두드러지며, 이에 더해 결함 어휘의
판정 기준 혼재라는 실무적 문제가 있다. 첫째, 분말 베드
세그멘테이션 모델은 여전히 수십 개의 어노테이션 레이어로 학습되는 장비별
지도학습 네트워크에 머물러 있으며, 장비와 소재 간에 시각 지식을 공유하는
메커니즘이 없다. 둘째, 평가가 도메인 내 일치도---배포된 모델의 자체 예측에
대한, 혹은 학습 분포에서 추출된 라벨에 대한 일치도---에 지배되어, 미확인
플랫폼으로의 일반화가 held-out 사람 ground truth를 기준으로 정량화되는
경우는 드물다. 셋째, 결함 맵이 물성 모델의 입력으로 점점 더 많이 사용되고
있음에도, 세그멘테이션 품질이 물성 예측에 대한 하류 가치로 평가되는 경우는
거의 없다. 이에 더해 데이터셋의 결함 어휘가 원인 기준과 외형 기준을
혼합하고 있어, 픽셀 단위 어노테이션의 판정 기준이 클래스마다 다르다.
DefectSegAM은 세 가지 결핍을 다루며, 어휘의 문제에 대해서는 실험과 별도로
외형 기준 분류 체계를 제안한다. 고정된 DINOv2 파운데이션
인코더는 시각 표현을 여덟 개 장비--소재 소스에 걸쳐 상각하면서 14.6~M
파라미터의 디코더만을 학습하고; 12{,}607개 양산 레이어에 대한 의사 라벨
사전학습은 배포된 DSCNN을 경쟁자에서 비용이 들지 않는 교사로 전환하며,
미세조정은 여섯 개의 L-PBF 소스와 두 개의 바인더 젯팅 소스에 걸친 125개의
사람 어노테이션 레이어를 취합하고; 평가 프로토콜은 학습에서 한 번도 보지
못한 장비의 이미지---그 수집과 라벨링이 진행 중이다---를 포함하는 held-out
사람 GT 벤치마크와, 세그멘테이션 정보원만을 유일한 변경 요소로 하는 하류
인장 특성 사례 연구를 결합한다. 마지막으로, 어휘 문제에 대한 응답인 외형 기준
결함 분류 체계는 픽셀 어노테이션의 판정 기준을 통일하기 위한 제안이며, 본
논문의 어떤 결과에도 사용되지 않았다.
